\section{Cartier selection on the leading field}
\label{sec:cartier}

Fix one germ, an independent assignment $b$, and the selected vertex
$w$ of Proposition~\ref{val:optimization}. Write
$\phi=\Phi(w)>0$ and retain the vectors $e,g_s,m_h$ of
\eqref{val:coefficient-vectors}. The constructions in this section
are made for these fixed characteristic-zero data. Reduction
characteristics will be arbitrarily large, but their required size
is not asserted to be uniform among germs or choices of tuples.

Fix any finite list of characteristic-zero regular eigenfunctions
$t_\nu$ whose $w$-initials are algebraic over the graded tuple
field of $f$, and any finite list of real cutoffs $T>0$.
In the reductions constructed below, $w$ denotes the valuation
obtained by keeping the parameter weights on a spread SNC chart
and restricting to the reduced germ. Write $w(a)=m_a\cdot b$
there; independence of $b$ makes the rational vector $m_a$ unique.
The finite tracked values and characters are preserved. We write
$w_{\rm orig}$ when distinguishing the characteristic-zero valuation.

\begin{proposition}\label{car:output}
There is a spread of the algebraic data, retaining the given weights
and cutoffs, such that, for all sufficiently large reductions of
characteristic $p$, there is a single product
$H$ of regular eigenfunctions with
\[
 \sigma=m_H=p\phi g_s+O(1).
\]
The error depends on the fixed characteristic-zero data but not on
$p$. For this $H$ and every vector of nonnegative integers
$(j_\nu)$, there are a regular nonzero eigenfunction $R$ and
integers $0\le k_i<p$ with
\begin{align}
 p m_R&=\sigma+\sum_\nu j_\nu m_{t_\nu}+k-(p-1)e,
 \label{car:C2-value}\\
 p\chi_R&=\chi_H+\sum_\nu j_\nu\chi_{t_\nu}
                   +\chi_{f^k}-(p-1)\chi_\eta.
 \label{car:C2-character}
\end{align}
Characters are written additively. For every preassigned cutoff
$T$, any collection of these reduced regular eigenfunctions
with distinct pairs \textup{(value, character)} and weights at
most $s(w_{\rm orig})T$ has size at most the characteristic-zero
number $\mathcal H(T)$ of Lemma~\ref{val:window}.
\end{proposition}

We first adjoin roots of the tuple and work on the residue field of
an extended valuation. Optimality gives a positivity inequality for
its residual form, and logarithmic Cartier turns that inequality into
a nonzero leading coefficient.
The same coefficient will work for every power vector $(j_\nu)$.

\subsection{The residue field and its form}

Choose a positive integer $N_0$ such that the coefficient lattice of
the value group is contained in $\frac1{N_0}\Z^r$. Adjoin roots
$x_j^{N_0}=f_j$, and extend $w$ to $\widetilde w$ on
$\widetilde L=L(x_1,\ldots,x_r)$. Its coefficient lattice is exactly
$\frac1{N_0}\Z^r$. Put
\begin{equation}\label{car:normalized-initial}
 K_0=\kappa(\widetilde w),\qquad q=n-r,\qquad
 a'=\operatorname{res}_{\widetilde w}
       \bigl(a x^{-N_0m_a}\bigr)\quad(a\in\widetilde L^*).
\end{equation}
The field $K_0$ is finitely generated of transcendence degree $q$
over $k_0$, and is generated by the $a'$ for regular eigenfunctions
$a\in L^*$.

Here are the field-theoretic details, which also account for the
possible failure of $w$ to be $G$-invariant. Let a finite abelian
group act on a field with a nontrivial real valuation, and let $D$
be its decomposition subgroup. Distinct orbit valuations are
inequivalent: a positive proportionality factor between valuations
in a finite orbit must be one. Given a $D$-invariant homogeneous
initial $\xi$ of value $\gamma$, take a lift $a$. By weak
approximation for the finitely many inequivalent rank-one valuations,
choose $h$ with
\[
 w(h-1)>0,\qquad u(h)>\gamma-u(a)
 \quad\text{at every other orbit valuation }u.
\]
Then $ha$ has initial $\xi$ at $w$ and value greater than $\gamma$
at all the other orbit valuations. Averaging $ha$ over $D$ preserves
both assertions, since $D$ permutes those other valuations. Taking
the sum over cosets of $D$ gives a group-invariant element with
initial $\xi$. Equivalently one uses the full group sum divided by
$|D|$. These strict inequalities work equally when $\gamma<0$.

Every homogeneous initial splits into $D$-eigenvectors. Characters
of a subgroup of a finite abelian group extend to the group. For
the Kummer extension $\widetilde L/L$, they are realized by
monomials in $x$; divide an eigen-initial by the corresponding
monomial initial and apply the invariant lifting just proved.
Thus the graded field upstairs is generated by those monomials
and initials from $L$. This proves that its value lattice is exactly
$\frac1{N_0}\Z^r$. Applying the same argument to $L/L^G$, rational
global eigenfunctions supply all subgroup characters. Their existence
follows, after quotienting by the kernel of the action, from a
normal basis for this finite Galois extension. A rational
eigenfunction is a quotient of regular eigenfunctions: multiply a
regular denominator by its group translates to obtain an invariant
regular denominator. Normalization by $x$ now proves the claimed
generation of $K_0$. The finite extension is still Abhyankar,
because graded independence survives finite extension. The
stratum description in Lemma~\ref{val:gauss-tools} identifies its
residue with a stratum function field, proving finite generation.

Choose valuation units $y_1,\ldots,y_q\in\widetilde L$ whose
residues $\bar y$ form a transcendence basis of $K_0/k_0$. Then
$x,y$ are a full Gauss tuple. Write the pulled-back form as
\begin{equation}\label{car:residual-form}
 \begin{aligned}
 \eta&=a\,d\log x_1\wedge\cdots\wedge d\log x_r
             \wedge d\log y_1\wedge\cdots\wedge d\log y_q,\\
 \rho&=a'\,d\log\bar y_1\wedge\cdots\wedge d\log\bar y_q.
 \end{aligned}
\end{equation}
By \eqref{val:log-frame}, $m_a=e$. The nonzero rational volume form
$\rho$ is independent of the chosen unit lifts and residual basis.
To check this, resolve both choices on a common proper smooth model.
At the generic centre stratum of codimension $r$, the vertical log
exponent matrix of $x$ is invertible. The residue of the horizontal
determinant changing the log frame is exactly the determinant
changing the restricted differentials on the stratum. The ratio
of the two frames is a unit with this residue, while normalization
by $x^{-N_0e}$ is the same in both. This proves the assertion and
permits resolving any additional finite list of rational supports.
If $q=0$, the wedge in $\bar y$ is $1$ and $\rho\in k_0^*$.

\subsection{Strict positivity over an affine residue model}

Choose a finitely generated $k_0$-algebra $A_0\subset K_0$, with
fraction field $K_0$, generated by normalized initials of regular
eigenfunctions. Include the normalized initials of every $h_i$ and
every detector. Let $J$ be the nonempty set of active generators,
and, for a divisorial valuation $z$ of $K_0/k_0$, put
\begin{equation}\label{car:l-definition}
 l(z)=-\min_{i\in J}\frac{z(h_i')}{v_i}.
\end{equation}

\begin{lemma}\label{car:strict-inequality}
For every divisorial $z$ one has
\begin{equation}\label{car:strict}
 A_\rho(z)+\phi l(z)\ge0.
\end{equation}
The inequality is strict when $z$ is nonnegative on $A_0$.
\end{lemma}

\begin{proof}
Both assertions are homogeneous in $z$, so we may take $z$ primitive.
Consider the lexicographic composite valuation $(\widetilde w,z)$,
whose second value on $a\in\widetilde L^*$ is $z(a')$. Adapt the
residual tuple so that $\bar y_1$ is a parameter for $z$ and the
other $\bar y_j$ are units with independent residues at its generic
centre. The tuple $x,y$ is Gauss for its two rows of lexicographic
weights. On its rational function field replace those rows by the
first row plus $t>0$ times the second row.

We require extensions which track exactly, and simultaneously, the
projected lexicographic values of finitely many elements. For each
test element, take its polynomial over the full tuple field. For
sufficiently small $t$, each coefficient value is the projection
of its lexicographic value. Pairwise equalities of term values give
a finite set of lexicographic candidates for the element's value.
Choose integer powers of the test elements so that their product
value separates all distinct tuples in the Cartesian product of
these finite candidate sets; include all candidates, whether or not
they are lowest values. Its actual lexicographic value is a Newton
root value of its minimal polynomial. The finitely many comparisons
determining the Newton polygon are preserved after projection for
all sufficiently small $t$. Hence the projected product value is
the value of a root over an algebraic closure of the projected
valued tuple field. This uses the Newton-polygon fact that a value
where the minimum of the term values is attained at least twice is
a root value for a polynomial with nonzero constant and leading
coefficients, as follows by factoring over a valued algebraic closure.
The embedding of the product field sending it to this root extends
to the whole finite extension $\widetilde L$. The values of all
test elements under that one embedding belong to their projected
candidate sets. Distinct separating product values remain distinct
for small $t$, so this forces the entire desired tuple of values.

Include the coefficient of $\eta$ in the adapted log frame, all
$h_i$, and every detector. The parameter $y_1$ already has prescribed
weight $t$ on the tuple field; it can also be included in the list.
The resulting extensions, and their restrictions $w_t$ to $L$,
are quasi-monomial. Indeed full graded transcendence degree is
preserved by finite extension and restriction, and the full-tuple
argument of Lemma~\ref{val:gauss-tools} applies. Discrepancies of
pulled-back forms agree with those on restriction. Exact tracking
therefore gives, for all sufficiently small $t>0$,
\begin{equation}\label{car:variation}
 A_\eta(w_t)=A_\eta(w)+tA_\rho(z),\qquad
 s(w_t)=s(w)-t l(z).
\end{equation}
In the second equality inactive generators remain inactive for
small $t$. In the first, the coefficient in the adapted residual
log frame has order $A_\rho(z)$. The $h_i$ retain positive values,
so $w_t$ stays centred; the $f_j=x_j^{N_0}$ retain their weights
$b_j$. Minimality of $\phi$ now gives \eqref{car:strict}.

Suppose equality holds and $z\ge0$ on $A_0$. Equations
\eqref{car:variation} imply $\Phi(w_t)=\phi$. Strict retraction
places every $w_t$ on the original cone complex. All detector
values are nondecreasing, since their normalized initials lie in
$A_0$. Their sum was maximal, so every detector value is constant.
The joint detector fibre is finite. Each of its valuations on $L$
has only finitely many extensions to the fixed finite field
$\widetilde L$. This contradicts the actual variation
$\widetilde w_t(y_1)=t$. Thus equality cannot hold over the affine
part. When $q=0$ there are no divisorial $z$, and both assertions
are vacuous.
\end{proof}

\subsection{A projective model with controlled top sections}

Choose a positive integer $M$ with $M/v_i\in\Z_{>0}$ for all
$i\in J$, and set $u_i=(h_i')^{M/v_i}$. Resolve the rational map
defined by these functions on a smooth projective model. They are
generating sections of a Cartier divisor $D_0$ in its chosen
rational trivialization, and, after pullback,
\begin{equation}\label{car:linear-system}
 z(D_0)=-\min_{i\in J}z(u_i)=M l(z).
\end{equation}
Neither the functions in this trivialization nor $D_0$ need be
effective.

We can choose such a smooth projective model $Z$, a reduced SNC
divisor $B_Z$, and an ample integral divisor $P$ supported on $B_Z$
with the following properties. The model maps to a normal
compactification of the normalization of $\Spec A_0$; the divisor
$B_Z$ contains the supports of $D_0$ and $\operatorname{div}(\rho)$;
and
\begin{align}
 P_T&\ge\max\{0,(D_0)_T\}
 &&\text{for every component }T\text{ over infinity},
 \label{car:infinity}\\
 \beta/\rho&\in A_0
 &&\text{for }\beta\in H^0(Z,\Omega_Z^q(\log B_Z)(P)),
 \label{car:top-sections}\\
 H^a\bigl(Z,\Omega_Z^i(\log B_Z)(P+jD_0)\bigr)&=0
 &&\text{for }a>0,\ 0\le i\le q,\ -|J|\le j\le0.
 \label{car:initial-vanishing}
\end{align}

We spell out the construction because negative coefficients of $P$
on the finite part are useful. Normalize a projective closure of
$\Spec A_0$, giving ample effective Cartier infinity. Let $Z_0$
be a smooth projective model resolving the map by the $u_i$, with
an effective ample divisor $P_0$. Choose a nonzero conductor
$c_0\in A_0$ from its affine normalization into $A_0$. Choose
$0\ne g\in A_0$ vanishing on every finite height-one prime of
that normalization lying under the supports of
$D_0,P_0,\operatorname{div}(\rho)$, or $\operatorname{div}(c_0)$.
Resolve these supports, those of $g$, and the pullback of infinity
by a sequence of blowups $\mu:Z\to Z_0$. Include the exceptional
divisors in $B_Z$ and retain the notation $D_0$ for its pullback.
Choose an effective exceptional divisor $F$ with $-F$ relatively
ample. For $m$ large, $m\mu^*P_0-F$ is ample. For $c\gg m$ and
then $h\gg c$, take a sufficiently large positive multiple of
\begin{equation}\label{car:P-construction}
 m\mu^*P_0-F-c\operatorname{div}(g)+hH_\infty.
\end{equation}
The principal term does not change ampleness, and the infinity
term is nef as a pullback of an ample divisor. Its coefficients
can dominate all required infinity coefficients. At the generic
points of finite height-one primes on the normal affine base the
birational map is an isomorphism. There the coefficients of $P$
are nonpositive and are as negative as needed at the finitely
many bad primes and finite primes in $B_Z$. A top log section
$\beta$ therefore has $\beta/(\rho c_0)$ regular in codimension
one on the affine normalization. Normality and the conductor
give \eqref{car:top-sections}. Finally, taking a large multiple
in \eqref{car:P-construction}, Serre vanishing gives all the
finitely many conditions \eqref{car:initial-vanishing}.
For $q=0$, take $Z$ to be a point; the divisors and supports are
empty, and the same assertions hold.

\subsection{Logarithmic Cartier nonvanishing}

Spread this finite collection of data over a finite-type integral
subring of $k_0$, and take geometric fibres in arbitrarily large
positive characteristics $p$. We keep the same notation on these
fibres. After shrinking the spread, smoothness, the relative SNC
condition, global generation by the $u_i$, ampleness, and the
finite cohomological vanishings persist. Spread also a basis of
$H^0(Z,\Omega_Z^q(\log B_Z)(P))$ and polynomial expressions for
its ratios to $\rho$ in the chosen generators $a_\lambda'$ of
$A_0$. Cohomology and base change make this a basis after
reduction. The parameter base can be chosen smooth over an open
of $\Spec\Z$. Its geometric points lift to length-two Witt
vectors by smoothness, so the reduced smooth SNC pair has a
length-two Witt lift.

The finite range of vanishings suffices for every nonnegative
twist. To see this explicitly, let
$E=\Omega_Z^i(\log B_Z)(P)$ and $s_J=|J|$. The exact generating
Koszul complex of the $u_i$ ends, after twisting, in
\[
 0\longrightarrow E((d-s_J)D_0)\longrightarrow\cdots
 \longrightarrow E((d-1)D_0)^{\oplus s_J}
 \longrightarrow E(dD_0)\longrightarrow0,
\]
with the intervening exterior-power multiplicities understood.
For $d\ge1$ every preceding shift lies between $-s_J$ and $d-1$.
Starting from \eqref{car:initial-vanishing}, induction makes them
all acyclic in positive degree. Break the resolution into short
exact sequences of successive images, beginning at the left.
The long exact sequences show that every image is acyclic.
The final kernel has $H^1=0$, so the last map is surjective on
global sections and the last sheaf is acyclic. Iteration gives
\begin{equation}\label{car:koszul-generation}
 H^a(Z,E(dD_0))=0\ (a>0,d\ge0),\qquad
 H^0(Z,E(dD_0))=\sum_{|\alpha|=d}u^\alpha H^0(Z,E).
\end{equation}
This concerns the globally generated line bundle
$\cO_Z(D_0)$; it makes no effectiveness assumption on its rational
trivialization.

Put
\begin{equation}\label{car:Hp}
 k_p=\lfloor p\phi/M\rfloor,\qquad H_p=k_pD_0+P,
 \qquad A_p=\lfloor H_p/p\rfloor.
\end{equation}
For all sufficiently large $p$,
\begin{equation}\label{car:rho-section}
 0\ne\rho\in H^0(Z,\Omega_Z^q(\log B_Z)(A_p)).
\end{equation}
At a boundary prime $T$, membership is equivalent to
$A_\rho(T)+(H_p)_T/p\ge0$, since $A_\rho(T)$ is integral.
Writing $a_p=p\phi/M-k_p\in[0,1)$, this expression is
\[
 A_\rho(T)+\frac\phi M(D_0)_T
       +\frac{P_T-a_p(D_0)_T}{p}.
\]
On the finite part the limiting expression is strictly positive
by Lemma~\ref{car:strict-inequality}; the fixed, possibly negative,
coefficient of $P$ is harmless for large $p$. At infinity the
limiting expression is nonnegative and the error is nonnegative
by \eqref{car:infinity}. Only finitely many coefficients are being
checked, and their orders have been preserved in the spread.

There is a global top section in pole twist $H_p$ whose Cartier
image on the function field is nonzero. Here is the cohomological
argument, including the effect of signs in $H_p$. Use relative
Frobenius and suppress the perfect-field twists of its target.
The meromorphic log de Rham complexes have an inclusion
\begin{equation}\label{car:complex-inclusion}
 \Omega_Z^\bullet(\log B_Z)(pA_p)
 \longrightarrow \Omega_Z^\bullet(\log B_Z)(H_p),
\end{equation}
with the ordinary meromorphic differential. After Frobenius
pushforward it is a quasi-isomorphism. In an etale chart adapted
to the log axes, write an axis coefficient of $H_p$ as
$h=pa+r$, $0\le r<p$, also when $h<0$. The smaller coefficient
module allows exponents at least $-pa$, and the larger one
allows exponents at least $-h$. The extra exponents have nonzero
residue modulo $p$. Split the Frobenius module by these residues.
A summand with nonzero residue in a log direction is contracted
by the log Euler field divided by that residue. The residue-zero
summands coincide. With several axes use any direction of nonzero
residue; ordinary coordinates do not affect the contraction.
Relative Frobenius commutes with etale charts, proving the claim.

By the projection formula, the pushed-forward smaller complex
is the Frobenius log de Rham complex tensored with $\cO_Z(A_p)$.
For $p>q$, logarithmic Deligne--Illusie splitting
\cite[Section~4.2.3]{DeligneIllusie}, applied to the smooth SNC length-two lift,
gives a surjection from its degree-$q$ hypercohomology to
$H^0(Z,\Omega_Z^q(\log B_Z)(A_p))$ by Cartier.
For clarity, the logarithmic splitting can be constructed in the
same way as the usual one. Local lifted Frobenius maps carrying
each axis to its $p$th power times a unit give divided pullbacks
of closed one-forms, including divided pullback on $d\log$ of
an axis. These lift inverse Cartier. On overlaps, divided
differences of the lifted maps give homotopies; for a log axis
use the ratio of the lifted pullbacks minus one, divided by $p$.
The ratio is a unit congruent to one. Difference and ratio
identities make the homotopies additive on triple overlaps.
They define the degree-one map in the Cech resolution of the
pushed-forward complex. Exterior products and antisymmetrization
then give the maps in all degrees, since $p>q$ makes the
factorials invertible. The coordinate Cartier formula identifies
their cohomology maps with inverse Cartier. Together with
Frobenius in degree zero they give the required derived splitting.
The pushed-forward form bundles are flat because relative
Frobenius of a smooth variety is finite flat. Tensoring the
splitting with $\cO_Z(A_p)$ requires no lift of this last divisor.

By \eqref{car:koszul-generation}, every term of the larger
complex in \eqref{car:complex-inclusion} has vanishing positive
sheaf cohomology. Its degree-$q$ hypercohomology is consequently
represented by global top sections. Choose a class mapping to
the nonzero section $\rho$ in \eqref{car:rho-section} and represent
it by such a top section. At the generic point the edge map is
ordinary field Cartier $C$, in the sense of Cartier's original
operation and its logarithmic coordinate formulation
\cite{Cartier1957}\cite[Theorem~7.2 and its proof]{Katz1970}. Since $\rho$ is nonzero generically,
the representative has nonzero field Cartier image.

Expand that representative using \eqref{car:koszul-generation}
and the fixed polynomial expressions for the top-section ratios.
Cartier is additive and satisfies
$C(c\omega)=c^{1/p}C(\omega)$ over the perfect constant field.
Thus at least one rational monomial term has nonzero image:
\begin{equation}\label{car:Up}
 U_p=u^\alpha\prod_\lambda(a_\lambda')^{I_\lambda},
 \qquad |\alpha|=k_p,\qquad C(U_p\rho)\ne0.
\end{equation}
The multi-index $I$ belongs to a fixed finite set. The individual
monomial need not separately be a global section; only its field
Cartier image is used. Replace $u_i$ by $h_i^{M/v_i}$ and
$a_\lambda'$ by $a_\lambda$. This gives a product $H$, depending on $p$, of
regular eigenfunctions on the reduced germ, with
\begin{equation}\label{car:C1}
 \sigma:=m_H=k_pM g_s+O(1)=p\phi g_s+O(1).
\end{equation}
The implied constants depend on the fixed characteristic-zero
data but not on $p$. In dimension $q=0$, $Z$ is a point and
Cartier in degree zero is inverse Frobenius on a nonzero scalar;
the same selection and formula apply.

\subsection{Spreading the initials and proving the output}

The test functions $t_\nu$ fixed at the start of the section have
normalized initials satisfying
\begin{equation}\label{car:constant-tests}
 t_\nu'\in k_0^*.
\end{equation}
Indeed after adjoining the initial roots $x$, they have degree
zero and are algebraic over the graded field of independent
variable weights with degree-zero field $k_0$. A homogeneous
relation makes them algebraic over $k_0$, which is algebraically
closed.

All these data can be spread simultaneously with the valuation
calculation. Resolve on $\widetilde L$ at $\widetilde w$ the
supports of $x,y$, of the normalized form coefficient, and of
normalized representatives for all initials used. Include the
$h_i,a_\lambda,t_\nu$, the detectors, and any further prescribed
nonzero functions, in particular the ideal generators below.
At the resulting codimension-$r$ stratum the parameter weights
are independent and generate $\frac1{N_0}\Z^r$. Hence the
exponent matrix of $x$ in these parameters is unimodular.
Normalized representatives of order vector zero are actual
units, and restriction computes their residues. Spread these
identities on open charts, the identification of the stratum
field with $k(Z)$, a separating residual tuple $\bar y$, and
the inclusions in the Kummer diagram. Retain geometrically
integral strata and models where they were so in characteristic
zero. After shrinking, the finite field extensions remain
separable. The same real parameter weights define monomial
valuations in every such reduced chart; their restriction to
the reduced germ is again denoted by $w$. The coefficient
lattice, the finite tracked values, and the normalized residues
are unchanged. In particular \eqref{car:constant-tests} remains
a nonzero constant identity, and the centred generators ensure
centre at the reduced closed point. Spread also the group action,
its characters with identified roots of unity, normality, and
the canonical generator on the smooth locus with complement of
codimension at least two. Exclude primes dividing $N_0|G|$.
Every construction here involves only finitely many data over a
finitely generated subring of $k_0$.

\begin{proof}[Proof of Proposition~\ref{car:output}]
Set $Q_p=H\prod_\nu t_\nu^{j_\nu}f^k$. Choose $k$ in the indicated
box so that
$m_{Q_p}+e\in p(\frac1{N_0}\Z^r)$. This is possible because
$p\nmid N_0$. Define
\[
 R=\frac{C(Q_p\eta)}{\eta}.
\]
Ordinary Cartier sends regular top forms to regular top forms
on the smooth locus. There $\eta$ is a canonical generator,
so $R$ is regular. Normality extends it across the complement
of codimension at least two. Once nonvanishing is known,
semilinearity and compatibility with the group action give
\eqref{car:C2-character}.

Compute the value upstairs, where Cartier commutes with
separable pullback, as follows also from its field $p$-basis
formula. The tuple $x,y$ is an absolute $p$-basis over the
perfect constants. Indeed its monomials with exponents in
$[0,p)$ are independent over $p$th powers already in the graded
field: use the basis of the value lattice for $x$ and the
separating residual $p$-basis for $\bar y$. There are $p^n$
of them, so they form a field basis. In the log frame of
\eqref{car:residual-form}, write the coefficient of $Q_p\eta$ as
\[
 c=\sum_{I,J}c_{IJ}^{p}x^I y^J,
 \qquad 0\le I_i,J_j<p.
\]
Lowest terms cannot cancel by that graded independence.
Since $m_c=m_{Q_p}+e$ lies in
$p(\frac1{N_0}\Z^r)$, lowest terms have $I=0$.
Normalize by $x^{-N_0m_c}$, a $p$th power. The residue times
$d\log\bar y$ is
\[
 U_p\prod_\nu(t_\nu')^{j_\nu}\rho.
\]
Its residual Cartier image is nonzero by \eqref{car:Up} and
\eqref{car:constant-tests}. Thus the term with $I=J=0$ occurs
at the lowest value, and $m_{c_{00}}=m_c/p$. Log-frame Cartier
extracts precisely $c_{00}$ times that frame. Division by
$\eta$ gives
$m_R=(m_{Q_p}+e)/p-e$, proving
\eqref{car:C2-value} and nonvanishing. The selected $H$ depends
on $p$ but is unchanged for every power vector: multiplication
by any product of the nonzero constants $t_\nu'$ cannot kill
its residual Cartier image. Thus no infinite list of products
needs to be spread.

For the final assertion, fix a cutoff $T$ in characteristic
zero. The ideal $I_T=(v>T)$ is primary to the maximal ideal,
since sufficiently high powers of its finite set of generators
have $v$-value greater than $T$. Choose finite generators of
$I_T$, identities expressing containment of a high maximal-ideal
power, and finite identities for $G$-stability. Spread these,
keeping any local denominators invertible. After shrinking,
the finite quotient has constant length $\mathcal H(T)$.
Track the chosen generators in the monomial chart just described.
By \eqref{val:comparison}, each has characteristic-zero value
strictly greater than $s(w_{\rm orig})T$, and its value is
unchanged on the reduction. Nonnegativity on the reduced local
ring therefore gives
\begin{equation}\label{car:cutoff-inclusion}
 I_{T,p}\subset (w>s(w_{\rm orig})T).
\end{equation}
There is no assertion that the whole valuation $v$ is spread.
If reduced eigenfunctions with distinct (value, character) pairs
of weights at most this cutoff had a dependence modulo $I_{T,p}$,
split it into characters using stability and $p\nmid|G|$.
Within a character the distinct lowest value cannot cancel,
contradicting \eqref{car:cutoff-inclusion}. Their number is
therefore at most the unchanged quotient length.
\end{proof}

The order of choices in Proposition~\ref{car:output} will be used
repeatedly: first fix the germ, its selected vertex, the finite
characteristic-zero test list, and all cutoffs needed for that
application; then spread those data and let $p$ be arbitrarily
large. The estimates obtained from Lemma~\ref{val:window} may
be uniform along the original sequence even though the required
reduction characteristic is not.
